\documentclass[twoside,12pt,titlepage,a4paper]{article}
\usepackage{url}
% kentHarvard requires natbib
\usepackage{natbib}
% add line numbers
\usepackage{lineno}
\usepackage{xcolor}
\linenumbers
\definecolor{periwinkle}{rgb}{0.8, 0.8, 1.0}
\renewcommand{\linenumberfont}{\normalfont\bfseries\small\color{periwinkle}}
\usepackage[pass]{geometry}
\usepackage{graphicx}
\renewcommand{\baselinestretch}{1.3}
\usepackage{todonotes}
\usepackage{enumitem}

\title{Checkmate+ COMP6030 report}
\author{
    Ainan Cawley\\
    \texttt{acc59@kent.ac.uk}
    \and
    James Ledger\\
    \texttt{jhl28@kent.ac.uk}
    \and
    Luke Williams\\
    \texttt{lw571@kent.ac.uk}
    \and
    elendil7\\
    
    \\ \vspace{10mm}
    \includegraphics[scale=0.6]{Kent_Comp_294_RGB} \\
    School of Computing \\
    University of Kent \\ \vspace{10mm} \\ Word Count: TODO
}
\begin{document}

\newgeometry{hmarginratio=1:1}    %% make layout symmetric
\maketitle
\restoregeometry              %% restore the layout

\begin{abstract}
   TODO: This should provide a concise summary of the paper. Typically, it will be between 100 and 200 words in length.
\end{abstract}


    \section{Introduction}
    \label{Introduction}
    
    Chess is a centuries old board game that has adapted well to the internet age. \textit{Chess.com}, the largest online chess platform, has over ten million games played on it every day \citep{chessdotcom-about}. 
    
    The online nature of these platforms enable features that don't exist when playing chess with a physical board and pieces. \textit{Lichess.org} introduced voice recognition to make moves in 2023 \citep{lichess-voice}, and \textit{Chess.com} offers alternative animations when moving or capturing pieces. However, there is little published evidence on whether such features improve user enjoyment.
    
    The goal of our project is two-fold. Firstly, to create a website for playing chess that implements three features: 
    \begin{itemize}[noitemsep]
    	\item Computer opponents that have human-like move times.
    	\item Alternative move and capture animations
    	\item Voice-command move input
    \end{itemize}
    Secondly, to investigate whether these features affect user enjoyment through surveying and conversational user evaluation sessions.
    
    \section{Collaboration methodology}
    \label{Collaboration methodology}

    TODO: basically how we worked together as a group to complete this project. Issue tracking, sprints/agile stuff. KEEP IT BRIEF.


    \section{Project architecture}
    \label{Project architecture}

    TODO: For James an overview of the structure of the project. Microservices, how they communicate (GRPC), springboot, GraphQL, STOMP as protocol to backend (instead of raw websocket), etc.

    Subsections for each service (brief)


    \section{The computer opponent}
    \label{The computer opponent}
    
    Online chess platforms often offer the option to play against a computer opponent (also called a ``bot''). The moves that a bot makes is determined by a program called a ``chess engine''. For \textit{Checkmate+}, we created our own chess engine, covered in section \ref{Chess engine}.
    
    Another aspect to the experience of playing against a bot is how long they take to make their moves. For \textit{Checkmate+}, we implement a feature that neither \textit{Chess.com} or \textit{Lichess.org} have: bots that have human-like move times. This feature will be covered in section \ref{Human-like move time model}.

    \subsection{Chess engine}
    \label{Chess engine}

    TODO: For Luke
    This is where you go into detail about what you have done. You will need to decide the titles
    for these sections yourself; they will depend on the content of the project. These sections
    should summarise the technical and scientific achievements of the project.
    Depending on the nature of your project, these sections may include: a comparison of
    different approaches that you considered, accounts of experimental work, mathematical
    analyses, specifications, top-level architectural diagrams, results obtained, problems
    encountered, workarounds, user evaluations, performance measures, testing regimes and
    results, comparisons between different approaches adopted, comparisons with existing work
    on similar problems.
    In particular, you should give a mixture of general discussion of your work and particular
    examples. Too much general discussion and the reader cannot easily get a handle on what
    you are doing; too many specific examples and the document fails to "tell a story".

    \subsection{Human-like move time model}
    \label{Human-like move time model}
    
    The time a computer opponent takes to reply to a user's moves affects the user experience. While it's usually best to minimise delays or loading times in software applications, a bot that replies near-instantly isn't ideal.
    
    \begin{quotation}
    	Paradoxically enough, moving as quickly as possible is not necessarily desirable. Users tend to get frustrated once they realize how little time their Mac really spends to crush them at low levels. In the interest of promoting harmonious Human - Machine relations, we enforce minimum response times \citep{apple-chess-engine}.
    \end{quotation}
    
    When creating Mac OS X's Chess application, the developers at Apple added a minimum delay of two seconds before the computer play its moves to avoid frustrating users. A more nuanced approach could improve on this. Human players do not take a uniform two seconds per move: some moves are obvious and are played quickly, while others warrant a longer think. A move time model that reflects this would feel more natural to play against. It could also increase user enjoyment, as a bot taking noticeably longer to reply may give the user the feeling that they have posed a genuine challenge, even if the bot is not truly ``stumped''.
    
    \subsubsection{Data}
    
    To train such a model, we used game data from the Lichess.org open database
    \citep{lichess-db-univ}. We used the January 2019 dataset, which contains millions of games along with each player's clock times in centisecond precision.
    
    The first 100,000 games in the dataset were filtered to only include those where a clock annotation exists for every move, and where at least four moves were played. Abandoned games and non-standard variants were excluded, as were casual (unrated) games on the basis that they're taken less seriously and may introduce noise. 55.72\% of games made it through this filtering.
    
    Moves in each game were converted into rows in our dataset for training the model. Each row contains the following features:
    
    \begin{itemize}[noitemsep]
    	\item \texttt{initial\_time}, the number of seconds on the clock when first starting the game. This feature is the same for all moves within the same game.
    	\item \texttt{increment}, the number of seconds added to each player's clock after they make a move. Like \texttt{initial\_time}, this feature remains the same for all positions within the same game.
    	\item \texttt{time\_remaining\_for\_this\_position}, the time in seconds remaining on the player's clock for them to make their move.
    	\item \texttt{time\_opponent\_took}, the time in seconds the opponent spent on their previous move.
    	\item \texttt{number\_of\_legal\_moves}, the number of legal moves available to the player in the current position.
    	\item \texttt{in\_check}, whether the player's king is currently in check.
    	\item \texttt{was\_capture}, whether the move just played by the opponent was a capture.
    	\item \texttt{move\_number}, the full move number of the current position. This number increases as a game progresses.
    	\item \texttt{time\_taken\_for\_this\_position}, the target feature. How many seconds the player took to make their move.
    \end{itemize}
    
    Note that since players' clocks don't tick down for their first move on \textit{Lichess.org}, the first 2 moves of every game aren't converted into rows in our training dataset so that we don't have any null values.
    
    In the end, our dataset for training move time models has 3,481,887 rows. This work can be seen in \texttt{/EngineMoveTimeModel/src/main/python/PGN to CSV.ipynb}.
    
    For training the linear regression model, two engineered features were included to help model non-linearity in existing features:
    
    \begin{itemize}
    	\item \texttt{time\_pressure}, defined as $\frac{1}{1 + t}$ where $t$ is \texttt{time\_remaining\_for\_this\_position}, grows sharply as a player's clock runs low. The idea being that players play faster when low on time to avoid losing via timing-out.
    	\item \texttt{opening\_unfamiliarity}, defined as $\frac{1}{1 + e^{6 - n}}$ where $n$ is \texttt{move\_number}. This is a sigmoid that increases from 0 to 1 approximately around the point at which players leave memorised chess territory. This turning point was found empirically through correlation with the target feature to be around move six.
    \end{itemize}
    
    For training a convolutional neural network model, a feature representing the state of the chess board was included. This feature was a 3-dimensional array of booleans of size (12, 8, 8) representing the 12 possible pieces (pawns, knights, bishops, rooks, queens, and kings in both colours) and their presence on the 8 by 8 chessboard.
    
    \subsubsection{Modelling}
    
    The dataset was split 80\%/10\%/10\% into training, validation, and test
    sets. As a baseline, a dummy regressor predicting the mean move time
    achieved an MSE (mean squared error) of 73.33 on the validation set.
    
    A linear regression model trained on the original features achieved a
    validation MSE of 56.81. The training data was then cleaned of
    outliers using the IQR (interquartile range) method, and the two engineered features added, bringing the validation MSE to 62.69. Although this is numerically higher than the unfiltered model, the outlier removal produced a
    healthier distribution of predicted move times. Before, there were a notable number of negative predicted move times, which was much improved after outlier removal. 
    
    The final model achieved a test set MSE of 64.40, corresponding to an RMSE (root mean squared error) of approximately 8.0 seconds.
    
    We also trained a model with a convolutional neural network to make use of the board state as an additional feature. The state of the board, a (12, 8, 8) boolean array was given to a convolutional neural network with 3 hidden layers, resulting in 512 outputs. These outputs were concatenated with the original features and passed through dense layers. This achieved a best validation MSE of 51.77 (RMSE $\approx$ 7.2 seconds) --- an improvement of roughly 0.8 seconds in typical prediction error over the linear regression model. However, given that human move times are inherently noisy and the
    improvement is unlikely to be perceptible to a user, this gain did not
    justify the added complexity: each training epoch took approximately 43
    minutes, and serving a model with 1.15 million parameters would have
    required a dedicated microservice.
    
    \subsubsection{Implementation}
    
    The linear regression model was therefore chosen for implementation.
    Given its simplicity, it runs inline within \texttt{GameServer}'s
    \texttt{MatchWebSocketController}, rather than as a separate
    microservice.
    
    \section{Front-end}
    \label{Front-end}

    TODO: For Ainan and James

    Could talk about our choice of only having TOTP auth app stuff.

    The frontend redesign.


    \section{Voice commands feature}
    \label{Voice commands feature}

    TODO: For elendil7
    This is where you go into detail about what you have done. You will need to decide the titles
    for these sections yourself; they will depend on the content of the project. These sections
    should summarise the technical and scientific achievements of the project.
    Depending on the nature of your project, these sections may include: a comparison of
    different approaches that you considered, accounts of experimental work, mathematical
    analyses, specifications, top-level architectural diagrams, results obtained, problems
    encountered, workarounds, user evaluations, performance measures, testing regimes and
    results, comparisons between different approaches adopted, comparisons with existing work
    on similar problems.
    In particular, you should give a mixture of general discussion of your work and particular
    examples. Too much general discussion and the reader cannot easily get a handle on what
    you are doing; too many specific examples and the document fails to "tell a story".

	\section{Feedback}
	\label{Feedback}
	
	TODO: what we get from users from the poster fair survey and chess club playtest
	
	\section{Future development}
	\label{Future development}

    \section{Conclusions}
    \label{Conclusions}
    

    \vskip 0.2in
    \bibliography{report}
    \bibliographystyle{kentHarvard}

\end{document}
